\section{Methodology}
\subsection{Datasets}\label{sec:data}
We use three English full-text news corpora: ISOT~\cite{ahmed2017isot} (real articles from Reuters, fake from sites flagged by fact-checkers; 2016--2017), the McIntire ``Fake or Real News'' set~\cite{mcintire2017fakeorreal}, and WELFake~\cite{verma2021welfake}, a collection merged from several sources. Each document is represented as title plus body. ISOT's \texttt{subject} and \texttt{date} fields are never used. LIAR~\cite{wang2017liar} and FakeNewsNet~\cite{shu2020fakenewsnet} were considered but excluded (short claims; crawling required).

\subsection{Audit and cleaning}
\emph{Normalization.} Two passes: lowercase, punctuation and whitespace collapse; and additionally removal of URLs, handles, a leading ``CITY (Reuters) --'' dateline, agency tags and the word ``Reuters'', so that the same story cleaned differently still matches (used for matching only).
\emph{Matching.} Exact hashes (128-bit) of text and of title+text; a 300-character prefix match; and near-duplicates by MinHash (128 permutations, 32 bands $\times$ 4 rows) over word 5-shingles on the full text and on the first 100 words, with candidates re-checked by the true Jaccard similarity (threshold 0.8) or containment ($\geq 0.9$). A sensitivity tier uses 3-shingles with Jaccard $\geq 0.5$.
\emph{Cleaning.} Documents with fewer than 20 words in the lightly normalized body are removed within each dataset. Exact duplicates are identified by the lightly normalized body hash; groups with conflicting labels are removed entirely, and the first row is retained in each remaining duplicate group. Near-duplicate clusters under the primary Jaccard/containment rule are kept but recorded, and the later splits are group-aware on these clusters. The length filter precedes title concatenation and differs from the final whitespace word count.
\emph{Label direction.} Documentation of WELFake is contradictory, so we verified it on Reuters-tagged articles (21,266 under label 0, 16 under label 1): label 0 is real and label 1 is fake. After this correction, exact-text label conflicts between datasets fell to one per pair.

\subsection{Analysis set and pre-specified thresholds}
Because every cleaned ISOT and McIntire document has a normalized match in WELFake (Section~\ref{sec:audit}), the builder removes WELFake documents flagged by an artifact-stripped body hash, a 300-character prefix, or primary near-duplicate matching against ISOT and McIntire, leaving a \emph{WELFake residual} (hereafter, the residual). The historical audit dry run also used lightly normalized and title-plus-body exact hashes and retained 17,161 documents; the experiment input retains 17,162. The extra fake-labelled row is an exact lightly normalized body match to an ISOT row, but both bodies become empty after artifact stripping, so the builder's stripped hashes do not match them. This is a construction-policy difference, not a change to the frozen experiment input. Before seeing the overlap results, we specified (without formal registration) the rule: keep a dataset if every class has at least 3,000 documents and the minority class is at least 40\% of the dataset. The rule narrowly failed (McIntire real: 2,964; residual minority share: 39.2\%). We proceeded and report this as a deviation; Section~\ref{sec:robust} shows a duplicate-removal rerun.

A subsequent provenance reconstruction from the available raw CSVs produced an analysis file byte-identical to the frozen experiment input (61,197 rows; SHA-256 \texttt{29e96d7b7fe60a94}\allowbreak\texttt{a2652c704f2bd4be}\allowbreak\texttt{9b5bb7e45f540eab}\allowbreak\texttt{eb19320276788ab5}). Original raw-source hashes were not recorded, so this reproduction does not prove historical raw-source snapshot identity. Both identifiers of the retained exact-match pair occur in the archived exclusion list; inspection of all saved paired-validation sample manifests found neither identifier in any selected training or evaluation set. This does not establish that all residual overlap has been eliminated.

\subsection{Interventions}
Each intervention deletes spans from the text and is applied to training and test documents alike.
\begin{itemize}
\item \textbf{S (source-marker masking):} regular expressions for URLs and broken URLs, datelines and agency tags, outlet names, image credits, ``share/tweet/print'' prompts, advertisement markers, posted-at stamps, ``Source:'' credit lines, email addresses, embedded-tweet attributions, and @/\# handles. The patterns were written after inspecting label-predictive tokens, and S therefore removes \emph{overt} markers, not all source information. S is split post hoc into W (datelines, agency tags, outlet names) and B (all other web boilerplate).
\item \textbf{L (length matching):} within each dataset, documents are binned into ten deciles of word count, and the same number of real and fake documents is sampled per bin, so length carries no class information within a dataset (approximately, at decile resolution).
\item \textbf{E (entities):} spans labelled PERSON, ORG, GPE, NORP, LOC, FAC, EVENT, PRODUCT, WORK\_OF\_ART, LAW or LANGUAGE by spaCy's \texttt{en\_core\_web\_sm}~\cite{honnibal2020spacy} are deleted.
\item \textbf{T (temporal expressions, exploratory):} weekdays, capitalized month names with optional day and year, four-digit years, numeric dates and clock times. T was added after inspecting residual cues.
\item \textbf{Random controls} $R_S$, $R_E$, $R_{SE}$ ($R_W$, $R_B$, $R_T$ for the exploratory parts): for every document, delete the same number of whole words as the real intervention removed, at random positions.
\item \textbf{Combined controls} $L{+}R_E$ and $L{+}R_{SE}$ (added after the first results): length matching followed by random deletion of as many words as E (respectively S+E) removes; they are the matched controls of L+E and S+L+E.
\end{itemize}
For TF-IDF and DistilBERT, the archived analysis uses the $2^3$ factorial over S, L, E (eight cells) with the three random controls (for DistilBERT, $R_S$ and $R_{SE}$ were added after an initial three-seed run); T and the S split are exploratory; the combined controls $L{+}R_E$ and $L{+}R_{SE}$ and the RoBERTa arm were added after the TF-IDF and DistilBERT results had been inspected (Table~\ref{tab:status}). The archived RoBERTa arm omits S+L and is not a complete factorial replication.

\subsection{Experimental protocol}
\emph{Universe.} Up to 4,000 documents per dataset and class are drawn once (22,005 documents). \emph{Splits.} For each seed, 20\% of near-duplicate clusters of each dataset are held out as the in-domain test pool; the cross-dataset test sample is drawn from all documents of the target. \emph{Sampling.} For each dataset and seed, the training, in-domain and cross-dataset sample sizes are fixed (about 1,495 per class for training, 400--500 per class in-domain, about 795 per class cross-dataset) and \emph{equal across all conditions}, with exactly 50\% fake. In the archived runs preserved in Appendix~\ref{sec:archived}, the sampling stream includes the condition identifier, so different conditions generally select different documents from the same pools. Comparisons are paired by seed and dataset pair, but not by individual document; they therefore include condition-specific sampling variation. For a given condition and seed, document selections are shared across models. \emph{Seeds.} Eight seeds change the split and the sample. The initial analysis used five seeds. After inspecting the five-seed results we extended it to eight seeds to obtain more stable estimates. The extension was applied uniformly to all conditions and both models, and no condition was added or dropped, so the eight-seed results are an expanded-sample analysis, not a replication specified in advance. The TF-IDF conclusions were unchanged, but some DistilBERT estimates moved (the specific effect of entity masking went from $-0.022$ \ci{-0.028}{-0.015} at five seeds to $-0.010$ \ci{-0.025}{0.005} at eight), and only the eight-seed values are reported. The DistilBERT controls $R_S$ and $R_{SE}$ were added after an initial three-seed run showed that random deletion itself changes transfer. Every archived result uses seeds 0--7; the new validation retains those eight identifiers across all three models; the duplicate-removal checks cover TF-IDF and DistilBERT only.

\emph{Document-paired validation.} The main results use a separate \texttt{paired-documents-v1} run on the same analysis-set bytes (SHA-256 verified against the frozen original input). Before these validation outcomes were examined, the focused design and seeds 0--7 were fixed. The choices were informed by the archived results and were not formally registered. For each seed and corpus, the paired sampler draws one training, in-domain and cross-dataset selection per length-matching setting. Every masking condition and random control within that setting reuses those exact article identities, including their labels and cluster memberships. The length-matched selection is distinct from the unmatched selection by design. The same manifests are shared across models. Condition-specific transformation and training streams remain reproducible; article pairing does not imply identical random deletion positions or training randomness across conditions.

TF-IDF covers C0, S, L, E, S+L, S+E, L+E, S+L+E, $R_S$, $R_E$, $R_{SE}$, L+$R_E$ and L+$R_{SE}$ (13 conditions). Each transformer covers C0, E, $R_E$, L, L+E, L+$R_E$, S+L+E and L+$R_{SE}$ (eight conditions). The fixed allocation was seeds 0--2, 3--5 and 6--7 across three Kaggle sessions; each transformer session distributed independent fits over two T4 GPUs. Every fit evaluates on its source and both target corpora. New and archived cells are never mixed. Per-seed manifests, protocol files, result hashes and observed software records accompany the results.

\emph{Matched controls for length matching.} L alone is a resampling intervention: it changes which documents are drawn, not their text, and it is applied at the same sample sizes and 50\% class balance as every other condition. Its comparison condition is C0, so L versus C0 estimates the effect of the length-bin constraint together with the resulting change in sampled documents; a word-deletion control would not be the matched comparison for it. When length matching is combined with masking, the matched comparison is length matching followed by random deletion of the same size ($L{+}R_E$, $L{+}R_{SE}$), which we added. Length matching uses class labels when selecting both training and test samples. It is a diagnostic evaluation on constructed distributions, rather than a deployment procedure for an unlabeled target dataset.

\begin{table}[t]
\caption{Status of each analysis component.}
\label{tab:status}
\centering\scriptsize\setlength{\tabcolsep}{3pt}
\newcolumntype{L}[1]{>{\raggedright\arraybackslash}p{#1}}
\adjustbox{max width=\linewidth}{\begin{tabular}{@{}L{2.5cm}L{1.8cm}L{3.8cm}@{}}
\toprule
Component & Status & Fixed before results?\\
\midrule
Dataset-selection thresholds & primary & yes, without formal registration\\
S, L, E factorial (8 cells) & primary & design yes; S patterns written after inspecting label-predictive tokens\\
Random controls $R_S$, $R_E$, $R_{SE}$ & primary & TF-IDF yes; DistilBERT $R_S$ and $R_{SE}$ added after an initial three-seed run\\
Seeds & primary & extended from five to eight after inspecting five-seed results, for all conditions and both models\\
S split (W, B) and T & exploratory & no; added after inspecting results\\
Combined controls $L{+}R_E$, $L{+}R_{SE}$; RoBERTa arm & added & no; added after inspecting TF-IDF and DistilBERT results\\
Six matched-control contrasts and Holm family (Appendix) & descriptive & no; defined after the exploratory results\\
Duplicate-removal rerun (TF-IDF, DistilBERT) & robustness & after the audit\\
Document-paired validation & validation & design fixed before new outcomes; informed by archived results, not registered\\
\bottomrule
\end{tabular}}
\end{table}

\subsection{Models and metrics}\label{sec:models}
\emph{TF-IDF + LR:} word 1--2-grams, minimum document frequency 2, maximum 0.95, sublinear term frequency, at most 200,000 features, logistic regression ($C{=}1$, liblinear)~\cite{pedregosa2011sklearn}. \emph{DistilBERT}~\cite{sanh2019distilbert}, a distilled version of BERT~\cite{devlin2019bert}, with the Transformers library~\cite{wolf2020transformers}: \texttt{distilbert-base-uncased}, first 256 tokens, 2 epochs, batch 16, AdamW with learning rate $3\times10^{-5}$, weight decay 0.01, 6\% warm-up, mixed precision on an NVIDIA Tesla T4 GPU. \emph{RoBERTa}~\cite{liu2019roberta}: \texttt{roberta-base} with the same recipe. Both transformers reuse the exact samples of the TF-IDF arm. The archived TF-IDF cells were run on the first author's system and the archived DistilBERT cells on a Google Colab T4 GPU (library versions of the latter were not recorded). The added conditions, the RoBERTa arm and the checks were run on Kaggle with two T4 GPUs, using a runner that parallelises the independent training jobs over the GPUs and derives every random stream exactly as the original runners; a seed-0 re-run of conditions C0 and L reproduced the archived DistilBERT results exactly (AUC, F1, accuracy and AP, difference 0.0 in all 18 cells).

\emph{Metrics.} ROC-AUC is the headline metric. F1 of the fake class is also reported but is less informative: with balanced classes, predicting ``fake'' for everything gives F1 of 0.667 and AUC of 0.5.

\subsection{Statistics}\label{sec:stats}
The unit of inference is the \emph{seed}. The six ordered pairs within a seed share training sets and are therefore not independent replicates, and the models share the same documents for a given condition and seed. For every comparison we compute a paired difference within each seed and ordered pair (for example, cross-dataset AUC under condition S minus under C0), average it with equal weight over the six pairs to obtain one value per seed, and summarize the eight seed values by their mean and a 95\% $t$-interval (seven degrees of freedom). We also report how many seeds agree in sign with the mean. These intervals describe resampling and training variability for these three datasets under this protocol; all seeds draw from one fixed document universe, so they say nothing about other datasets, and the pooled average is a summary of six pairs that differ, which is why every pooled result is accompanied by per-pair results. Intervals and seed agreement are the basis for claims. For the six matched-control contrasts we additionally report exact paired sign-flip $p$-values with Holm adjustment~\cite{holm1979simple} (Appendix); they describe stability under resampling within one fixed document universe, the family was defined after earlier exploratory results had been seen, and no other comparison is corrected for multiplicity. Factorial effects use $\pm1$ coding: an effect is the average change when a factor (or interaction) goes from $-1$ to $+1$.

\emph{Comparing models.} To compare models we do not contrast ``significant'' with ``not significant''; we compute, for each seed, the difference between the effects of two models for the same intervention (for example $D_s=\text{effect}^{\text{TF-IDF}}_s-\text{effect}^{\text{DistilBERT}}_s$), which is paired because all arms use identical samples, and examine the eight values $D_s$. The archived full design contains eight factorial cells, three random controls and two combined controls. The paired transformer subset is specified above. The six-contrast family is retained in paired analyses; omitted transformer contrasts are assigned $p=1$ for Holm adjustment. All per-pair breakdowns and model-difference summaries are exploratory.
